\section{模型评价与结论}

\subsection{四问结果与模型特点}

四个问题在相同物理状态方程与结算定义下递进。问题一的确定性LP解释储能如何在容量和效率约束下转移能量；问题二由五倍补购机制推导采购倾向，再用残差场景处理不确定性；问题三区分允许修订的交易时点和每10分钟的储能动作；问题四将同日价格残差与净负荷残差联合使用。汇总如下。

\begin{table}[htbp]\centering\small
\caption{最终分问配置的334天开发期闭环评价结果}
\label{tab:annual_summary}
\begin{tabularx}{\linewidth}{LRRRR}\toprule
指标 & 问题2 & 问题3 & 问题4日前版 & 问题4滚动版\\\midrule
计划费用/万元 & 1350.80 & 1350.31 & 1428.01 & 1426.52\\
调整费用/万元 & 0.00 & 24.87 & 0.00 & 28.31\\
紧急费用/万元 & 50.48 & 11.97 & 51.40 & 12.19\\
实际总费用/万元 & 1401.28 & 1387.15 & 1479.41 & 1467.02\\
紧急电量/kWh & 85017.69 & 21029.90 & 80544.00 & 21293.41\\
紧急购电天数 & 113 & 119 & 106 & 116\\
紧急购电连续时段数 & 205 & 204 & 195 & 193\\
日总费用CVaR95/元 & 70731.74 & 65552.76 & 83044.31 & 77575.87\\
最高日费用/元 & 89306.82 & 67941.61 & 108080.37 & 81797.18\\
2月1日期初库存/kWh & 8875.19 & 8814.60 & 9251.54 & 9030.40\\
12月31日期末库存/kWh & 6240.32 & 6240.32 & 6269.86 & 6274.99\\
\bottomrule\end{tabularx}
\par\smallskip\parbox{\linewidth}{\footnotesize 问题2与问题4日前版沿用母模型；问题3与问题4滚动版在零时及日内均启用候选筛选。预测器与实时执行器沿用原冻结配置，筛选适用问题与时点依据开发期比较确定。实际费用不含候选评分的终端偏离项；CVaR95为实际日总账单的经验尾部均值。}
\end{table}


问题一费用为35126.95元，相对无储能节费26.90\%；问题二、三评价期主费用分别为1401.28万元、1397.58万元，动态电价两案分别1479.41万元、1476.92万元。机制推导说明采购偏高可以源于经济最优的非对称误差代价，不应自动等同于风险厌恶。固定合同的执行回放则把规划偏差与真实账单分开，避免将更接近参考轨迹直接当成更省钱。

主结果逐段满足供需平衡、储电上下界、功率限制、充放电互斥及日际连续性。未来实际值与未来发布预报的扰动没有改变已执行前缀。所有费用与时段分组由同一结果源生成，保留完整源代码、原模板及修订前版本，便于追溯图表中的每个数值。

\subsection{适用边界与可继续改进之处}

第一，只有一个年度且一月校准窗口较短，修订已经参考过旧版回测结果，不能声称跨年份泛化或从未接触的独立测试。第二，预测与控制采用有限候选的坐标搜索，28场景、额外风险与备用关闭等设置预先固定，并未穷举全部可行策略。第三，上层共享储能参考和实际控制之间仍有差异，零时尚未精确计入以后1.5倍修订的机会，也未求解完整跨日随机价值。第四，终端惩罚是辅助近似，比较时仍须报告末库存；模型未加入赛题以外的网络限额、设备寿命退化和交易约束。

后续改进应首先对齐上层场景评估与下层因果控制，再以合法预报发布时点构造规模可控的场景树。风险、备用、内层LP与跨日续接在不同设置下可能产生相反的费用变化，采用依据应是目标明确、信息合法且可以复现的结算收益。
